% ==============================================================================
% capitulos/01-introducción.tex - Introducción del Proyecto
% ==============================================================================

\chapter{Introducción}
\begin{flushleft}
\justifying
El presente documento lleva el avance realizado del TT 2026-B009 realizado en el semestre 2025-2, este trabajo se realizo con la Defensoría de Derechos Politécnicos (DDP) del Instituto Politécnico Nacional (IPN), con el objetivo de identificar áreas de oportunidad y proponer una solución tecnológica que optimice sus procesos internos.

El Instituto Politecnico Nacional es una institución educativa publica mexicana que se ha consolidado como una de las principales casas de estudio en el país, ofreciendo una amplia gama de programas académicos en diversas áreas del conocimiento \cite{1}.

La Defensoría de Derechos Politécnicos (DDP) del Instituto es un órgano autónomo que actuará con independencia de las autoridades del mismo y que tendrá por objeto llevar a cabo la promoción, protección, defensa, estudio y divulgación de los derechos de los miembros de la comunidad politécnica \cite{3}.

La DDP ha presentado inconvenientes en sus procesos de trabajo. Estos procesos de trabajo están establecidos en el Acuerdo de Creación de la Gaceta Politécnica No. 622 \cite{3}, la cual dicta que será competente para conocer de las quejas relacionadas con presuntas violaciones a los derechos de los miembros de la comunidad politécnica cuando éstas sean imputadas a las autoridades administrativas o académicas del Instituto, en relación con los actos u omisiones en que incurran en la aplicación de la normatividad del Instituto.

Una queja es el documento principal que da pie a una serie de procesos con el fin de salvaguardar los Derechos Politécnicos, Si bien, la DDP rige su funcionamiento a través del Manual de Organización \cite{4} y el Manual de Procedimientos \cite{5}, existen situaciones internas que aplazan la solución:
\begin{itemize}
\item Falta de personal.
\item Acumulación de trabajo
\item Duplicidad de quejas.
\item Falta de comunicación interna 
\end{itemize}

Por ende, en el desarrollo de este documento, se presenta la propuesta solución, con el fin de mejorar las áreas de oportunidad identificadas, esta propuesta denota el desarrollo de un sistema web que permita el intercambio de información, a su vez, permitirá a la persona que presento la queja (quejoso) observar el proceso de la misma.

\section{Planteamiento del Problema}
\end{flushleft}
\justifying
Una vez, establecida que es la DDP y las situaciones que provocan lentitud en los resultados, desarrollaremos el problema a resolver explícitamente.

\subsection{Contexto y situación problemática}
En primer término, se identifica una vulnerabilidad crítica en la estructura operativa de la Defensoría: la desproporción entre la demanda de atención a la comunidad Politécnica y la capacidad del capital humano especializado. Actualmente, la función principal de resolución de quejas recae de forma exclusiva sobre un cuerpo técnico de cinco abogados, quienes tienen la responsabilidad de gestionar la totalidad de las quejas declaradas procedentes dentro de una comunidad que abarca cientos de miles de integrantes.

Bajo el modelo de gestión vigente, este personal de alta especialización jurídica se ve absorbido por una densidad administrativa abrumadora. Al carecer de una plataforma centralizada, los abogados deben dedicar una fracción significativa de su jornada a tareas procedimentales, tales como el registro manual de expedientes, el seguimiento analógico de plazos procesales y la organización descentralizada de evidencias. Esta carga de trabajo actúa como un factor de perdida del tiempo efectivo, limitando el análisis profundo que cada caso amerita para garantizar el debido proceso.

\subsection{Problema a resolver}
Entonces, buscamos mejorar la optimización del proceso de gestión y resolución de quejas, con el fin de reducir los tiempos de respuesta y garantizar una atención oportuna al quejoso. Actualmente, la eficiencia se ve comprometida por la gestión manual de las etapas procesales.

El flujo de trabajo que se busca intervenir se compone de las siguientes fases críticas:

\begin{enumerate}
\item \textbf{Recepción de queja:} Captura inicial y registro del expediente.
\item \textbf{Análisis de competencia:} Evaluación jurídica para determinar la procedencia.
\item \textbf{Investigación:} Recopilación de pruebas y diligencias administrativas.
\item \textbf{Resolución:} Emisión del dictamen o recomendación correspondiente.
\item \textbf{Cierre:} Notificación a las partes y archivo del caso.
\end{enumerate}

Para dar respuesta a esta problemática, en el siguiente capítulo se describe la propuesta de solución, la cual busca integrar las capacidades tecnológicas necesarias para transformar el modelo de gestión actual.

\section{Solución Propuesta}
Se propone el desarrollo de una plataforma web diseñada para centralizar, automatizar y optimizar el ciclo de vida de las quejas dentro de la DDP. La solución se fundamenta en una arquitectura basada en microservicios y contempla la implementación de un módulo de asistencia inteligente en la etapa de primer contacto, el cual utiliza técnicas de PLN para asistir al área de Primer Contacto y determinar la competencia de la queja. Se permitirá la gestión basada en roles, garantizando que cada actor interactúe con la información de acuerdo con sus facultades normativas, asegurando el seguimiento total del expediente desde su recepción hasta su resolución final.

\section{Justificación}
La implementación de esta plataforma surge de la necesidad imperativa de optimizar los procesos de la DDP, entendiendo que este órgano representa el último recurso de protección para la comunidad. El retraso en la resolución de los casos no es solo un fallo administrativo, sino que puede derivar en la desprotección efectiva de los derechos de la comunidad Politécnica.

Al tener la información en un entorno web, se busca mitigar el impacto de la desproporción actual entre la alta demanda de usuarios y el limitado cuerpo técnico de cinco abogados. Esta solución permitirá:

\begin{itemize}
\item Erradicar la dispersión de documentos mediante un repositorio digital único.
\item Eliminar la duplicidad de quejas, optimizando el uso del tiempo del personal especializado.
\item Reducir significativamente los tiempos de respuesta, mejorando la gestión integral y garantizando una atención oportuna y eficiente.
\end{itemize}


\section{Objetivos}

\subsection{Objetivo General}
Desarrollar una plataforma web basada en una arquitectura de microservicios mediante la integración de módulos de gestión documental y un motor de procesamiento de lenguaje natural para sistematizar el flujo operativo de recepción, clasificación y seguimiento de quejas de la DDP.

\subsection{Objetivos específicos}
\begin{itemize}
    \item \textbf{Analizar} el marco normativo de la Defensoría de los Derechos Politécnicos y sus manuales de procedimientos para definir las reglas de negocio que permitan la transición de la gestión analógica a un flujo de trabajo digitalizado.
    
    \item \textbf{Diseñar} la arquitectura del sistema basada en microservicios, asegurando la integridad de los datos y una experiencia de usuario (UX) accesible mediante:
    \begin{itemize}
        \item \textit{Arquitectura:} Establecer un diseño modular que soporte la alta demanda de la comunidad Politécnica y garantice la seguridad de la información sensible.
        \item \textit{Base de datos:} Estructurar un esquema centralizado que elimine la duplicidad de registros y permita la trazabilidad total del expediente.
        \item \textit{UI/UX:} Proyectar una interfaz intuitiva para el quejoso y un panel de gestión eficiente para el cuerpo técnico jurídico.
    \end{itemize}

    \item \textbf{Implementar} un módulo mediante técnicas de Procesamiento de Lenguaje Natural (PLN) para asistir en la clasificación preliminar de las quejas y determinar su competencia institucional.

    \item \textbf{Desarrollar} los componentes centrales de la plataforma web, garantizando la operatividad de las siguientes funciones:
    \begin{itemize}
        \item \textit{Gestión de expedientes:} Digitalización del ciclo de vida de la queja (recepción, investigación y resolución) para optimizar el tiempo efectivo del personal jurídico.
        \item \textit{Seguimiento:} Establecer un sistema de consulta para el quejoso que brinde transparencia sobre el estatus de su trámite.
        \item \textit{Gestión por roles:} Implementar controles de acceso estrictos basados en las facultades normativas de cada actor dentro de la Defensoría.
    \end{itemize}
\end{itemize}

\section{Alcances y Limitaciones}
Tras definir los objetivos, se plantean claramente los alcances y limitaciones de la propuesta para contextualizar su implementación y las expectativas pensadas para los usuarios finales.

\subsection{Alcances}
Los alcances describen las funciones y procesos específicos que el sistema cubrirá de principio a fin, delimitando las capacidades técnicas que se entregarán al finalizar el proyecto.

La plataforma web permitirá el registro, seguimiento y gestión integral de las quejas recibidas por la DDP, abarcando desde el primer contacto del ciudadano ya sea por vía digital, telefónica o presencial hasta la resolución final del expediente, incluyendo notificaciones automáticas y generación de reportes estadísticos. 

El sistema implementará un clasificador basado en Procesamiento de Lenguaje Natural (PLN) para asistir en la determinación automática de competencia de las solicitudes, categorizándolas por materias como derechos humanos, amparo, penal o civil, con una interfaz intuitiva que sugiera asignaciones preliminares y reduzca tiempos de procesamiento en hasta un 40\% según pruebas piloto. 

Asimismo, contará con un repositorio digital seguro para el almacenamiento de documentos y evidencias, cumpliendo estándares de encriptación AES-256 y autenticación multifactor, facilitando el acceso compartido y la trazabilidad de versiones para auditorías internas. 

Se desarrollarán módulos específicos y personalizados para cada rol institucional —Abogado Asesor, Subdefensor, Recepcionista, Área de Primer Contacto y Titular—, con dashboards adaptados que incluyan alertas prioritarias, flujos de trabajo colaborativos y herramientas de firma electrónica avanzada. 

Finalmente, se proporcionará un sistema de consulta pública mediante folios digitales únicos, permitiendo a los usuarios externos rastrear el estatus de sus quejas de forma anónima o autenticada, fomentando la confianza ciudadana en la institución.
\subsection{Limitaciones}
Las limitaciones definen hasta donde llega el sistema, especificando aquello que el software no puede realizar por sí solo o que depende de factores externos y criterios humanos.

El sistema no sustituirá el criterio jurídico de los abogados de la DDP, sino que fungirá como una herramienta de apoyo en la clasificación inicial de quejas, permitiendo una triaje eficiente que acelere el flujo de trabajo sin comprometer la autonomía profesional. 

La plataforma no realizará resoluciones automáticas de quejas ni emitirá dictámenes legales vinculantes, ya que su diseño prioriza la revisión humana en todas las etapas críticas para garantizar el cumplimiento de normativas éticas y legales, como la protección de datos personales bajo la Ley Federal de Protección de Datos Personales en Posesión de los Particulares (LFPDPPP). 

Asimismo, el clasificador basado en IA requerirá un período de entrenamiento y ajuste continuo con datos reales anonimizados de casos históricos, lo que podría tomar varios meses para alcanzar niveles óptimos de precisión (por encima del 90\% en pruebas preliminares). En etapas iniciales, operará de manera asistida con supervisión humana obligatoria, donde los abogados validarán las clasificaciones sugeridas para mitigar sesgos algorítmicos, falsos positivos o negativos, y adaptarse progresivamente a la variabilidad del lenguaje jurídico en contextos mexicanos.
 
Adicionalmente, la efectividad del sistema dependerá de la calidad y volumen de los datos de entrenamiento, por lo que se recomienda un protocolo de retroalimentación iterativa para su mejora continua.